\section{Complete clock fibres and finite spatial records}
\label{sec:v84-finite}
\subsection{The entire fibre below the sharp polynomial threshold}
The implicit-function lower bound has a global form. This form also
specifies which additional information removes the remaining ambiguity.

\begin{theorem}[Global polynomial ambiguity fibres]\label{thm:v84-fibre}
Use $J=q+2$ clocks, let $I\subset(-\infty,T_1)$ be an open interval,
and put
\[
 w(s)=\prod_{j=1}^{q+2}(T_j-s)^{-1},\quad F'=w,\quad
 d=(-1)^{q+2}[T_1,\ldots,T_{q+2}]f.
\]
For $d\ne0$, the complete fibre of
$f_j=P(T_j)+\log(T_j-x)-\log(T_j-y)$, with $P\in\cP_q$ and
$x,y\in I$, is parametrized without repetitions by
\begin{equation}\label{eq:v84-fibre}
 u\in U_d:=\{u:u,u+d\in F(I)\},\qquad
 y(u)=F^{-1}(u),\quad x(u)=F^{-1}(u+d).
\end{equation}
For each such $u$ there is exactly one interpolating nuisance polynomial
$P_u$. Every nonempty fibre is a connected smooth open interval. Its
orientation is fixed by the sign of $d$. Knowing either action selects
at most one point of the fibre. If $d=0$, there is no unequal-action
solution; allowing equal actions gives the diagonal action fibre and
the uniquely determined sampled polynomial.
\end{theorem}
\begin{proof}
The unique annihilator of the degree-$q$ polynomial evaluations gives
$d=\int_y^xw=F(x)-F(y)$. Since $F$ is strictly increasing, this is
equivalent to \eqref{eq:v84-fibre}. The set $U_d$ is the intersection
of two open intervals and is an open interval whenever nonempty.
For every parameter $u$, subtracting the corresponding logarithms
from the sampled $f$ gives a vector killed by that annihilator. It
therefore lies in the $(q+1)$-dimensional polynomial evaluation space
and determines exactly one $P_u$. Conversely every solution has this
form. Smoothness follows from $w>0$ and the fixed interpolation matrix;
$y'(u)=1/w(y(u))>0$ proves regularity and absence of repetitions.
If $d=0$, strict monotonicity gives $x=y$.
\end{proof}
For the full interval $I=(-\infty,T_1)$ one may take the explicit
primitive
\begin{equation}\label{eq:v84-primitive}
 F(s)=-\sum_{i=1}^{q+2}
 \frac{\log(T_i-s)}{\prod_{h\ne i}(T_h-T_i)}.
\end{equation}
Partial fractions show that $F'=w$, $F(-\infty)=0$ and
$F(T_1-)=+\infty$. Thus $U_d=(\max(0,-d),\infty)$.
For $q=0$, with $\Delta=T_2-T_1$, the parametrization is elementary:
\[
 F^{-1}(u)=T_1-\frac{\Delta}{e^{\Delta u}-1}.
\]
For larger $q$, \eqref{eq:v84-primitive} and a single monotone inverse
specify the whole fibre, rather than just a neighbourhood supplied by
the implicit function theorem. Additional parameter bounds may cut
this interval; connectedness is asserted for the declared unrestricted
polynomial nuisance class, not for every additional constrained class.

\begin{corollary}[Finite-record consequence of an exact fibre]
\label{cor:v84-fibre-risk}
Fix positive full-rank two-component channels and choose two distinct
points of a fibre whose ordered action pairs have distance $D>0$.
They give identical projective matrices, and therefore identical laws
for any number of independent records drawn from those matrices.
Every estimator of the channel-labelled action pair has worst-case
expected Euclidean error at least $D/2$ on these two models, and
worst-case mean squared error at least $D^2/4$. Two points may be
chosen within any sufficiently small strict-margin fibre segment.
\end{corollary}
\begin{proof}
The fixed log probability ratios determine the two normalized component
weights and hence the matrices. Under their common record law, the
triangle inequality gives the first bound after taking expectations;
the parallelogram identity gives the squared bound. The scalar fibre
varies the actions, not merely the common detector gauge. Decreasing
a common retention baseline gives a strict upper retention bound on
the finite clock interval, and continuity preserves strict constraints
on a small segment.
\end{proof}
On the identifiable side, a finite design satisfying
\eqref{eq:v84-design} has a deterministic error guarantee. If both a
true scalar model and a fitted model approximate the observed samples
within $\varepsilon$, their nuisance and action distance is at most
$2\varepsilon/c_K$. With measurement error $\varepsilon$ and detector
misspecification bounded by $\eta$ in the sampled supremum norm, the
same argument gives $2(\varepsilon+\eta)/c_K$ for any fit with that
residual tolerance. This statement uses a specified identifiable
compact class; it does not identify an unrestricted detector error.

\subsection{Pointwise compact classes}
For the following theorem take the polynomial model with $J=q+3$.
Fix $q,u,v$, residual bounds $r,R$, centred coefficient bounds and
probability, rank, spectral, action and column-signature margins as in
\eqref{eq:v83-margins}--\eqref{eq:v83-signature-gap}. For every
$2\le B\le\min(u,v)$ these inequalities define a compact pointwise
parameter class $\cK_B$, after bounding all actions and centred
coefficients. The channels lie in their probability simplices. Use
the Euclidean metric on the centred parameters modulo permutation.
Empty classes are omitted.

There is a constant $C_0$ such that, for any two exact normalized
matrix tuples in the same $\cK_B$, the quotient distance of their
parameters is at most $C_0$ times their matrix distance. Indeed the
exact inverse is injective modulo a finite permutation group and has
a smooth local extension at every point, by
Theorem~\ref{thm:v83-main} and Lemma~\ref{lem:v83-local}. A finite
subcover gives the bound for nearby parameter pairs. Pairs outside
these neighbourhoods have a positive minimum observed separation by
compactness and injectivity, which supplies the remaining bound.
Taking the maximum over the finitely many $B$ gives $C_0$.
This constant may depend on the separation modulus of this entire
compact pointwise class; it is not estimated by an unverified local
Jacobian alone.

For a noisy tuple, first threshold the rank as below and then minimize
$(\max_j\|P_j^{\rm model}-\widehat P_j\|_F)^2$ over $\cK_B$. A minimizer exists by
compactness. A fixed measurable tie-breaking rule may be used. If the
true tuple is in $\cK_B$ and the data error is $\eta$, the fitted and
true observation tuples differ by at most $2\eta$; the fitted centred
parameter error is consequently at most $2C_0\eta$ modulo permutation.
This is an existence construction, not a claim of a polynomial-time
optimization algorithm. Crucially it uses the specified class, not
a supplied reference covering or a reference set of component names.

\begin{theorem}[Finite records recover the covering and its monodromy]
\label{thm:v84-mesh}
Let $Q$ have a fixed finite triangulation of mesh diameter at most $h$,
with $N_h$ vertices. Distances and simplex coordinates are fixed.
Assume the exact field takes values in the classes $\cK_B$ above,
with one unknown constant rank $B$. Along lifts of each simplex, assume
that the channel columns and the full centred field have Lipschitz
constants at most $L_U$ and $L_\Xi$, respectively, uniformly over
all local enumerations. These are prior regularity bounds, not
quantities certified from arbitrarily sparse data.

At each vertex and clock observe $n$ independent readout pairs with
law $P_j$ at that vertex. Put
\[
 \eta_n=\sqrt{\frac{uv}{2n}
           \log\frac{2JuvN_h}{\zeta}},\qquad e_n=2C_0\eta_n,
                    \quad 0<\zeta<1.
\]
If
\begin{equation}\label{eq:v84-mesh-margins}
 \eta_n<\sigma_*/4,\qquad 2L_Uh+4e_n<s_* ,
\end{equation}
there is a reconstruction from these finite records which, with
probability at least $1-\zeta$, recovers $B$, the component covering
up to isomorphism over $Q$, and its monodromy. It also gives a
continuous piecewise affine centred field on that covering satisfying
\begin{equation}\label{eq:v84-mesh-error}
 d_0([\widehat\Xi],[\Xi])\le e_n+L_\Xi h.
\end{equation}
The guarantee is uniform over the stated classes and all covering
types satisfying these bounds. The rank threshold uses the supplied
signal lower bound $\sigma_*$, not unrestricted model-order discovery.
No $C^m$, $m\ge1$, conclusion is asserted for piecewise affine output.
\end{theorem}
\begin{proof}
For each empirical matrix entry, the bounded-indicator exponential
inequality \cite{V84Hoeffding} gives error at most
$\sqrt{\log(2JuvN_h/\zeta)/(2n)}$ simultaneously at all vertices,
clocks and entries, except on an event of probability $\zeta$.
The Frobenius error of every matrix is then at most $\eta_n$.
Singular-value thresholding at $\sigma_*/2$ recovers $B$ at every
vertex: the nonzero exact singular values exceed $\sigma_*$ and the
remaining ones vanish. If the observed ranks are inconsistent, the
procedure may reject; this does not occur on the stated event.
Use the compact-class fit at each vertex and enumerate its columns
arbitrarily. Its component error is at most $e_n$ after some matching
to the true fibre.

On an edge, transport a true component along that edge. Its true
channel moves by at most $L_Uh$. The fitted columns for its two ends
are thus at distance at most $L_Uh+2e_n$. The distance to a fitted
column of a different component at the other end is at least
$s_*-L_Uh-2e_n$. Condition \eqref{eq:v84-mesh-margins} separates these
two bounds strictly. Nearest-column matching therefore gives exactly
the true edge transport as a unique bijection, without knowing that
transport in advance.

Every triangular face is contractible in $Q$, so the true transports
multiply to the identity around it. The recovered edge permutations
therefore satisfy the face relations exactly. Gluing copies of the
simplices with these permutations constructs a covering isomorphic to
the true covering. Monodromy along any edge loop is its permutation
product; the face relations and cellular approximation give the same
representation for every loop in $Q$. This also handles disconnected
total covering spaces. A failed face relation may be declared a
rejection outside the good event; a small numerical mismatch is not
silently substituted for an exact permutation cocycle.

Within each simplex, transport the fitted vertex labels to one of its
fibres and interpolate by barycentric coordinates. The edge relations
make these definitions agree on common faces. The centred gauge is
preserved by this affine interpolation. In the corresponding true
trivialization the noise contribution is bounded by $e_n$ and the
variation from a vertex by $L_\Xi h$. Their sum proves
\eqref{eq:v84-mesh-error} on the reconstructed covering and hence in
the quotient metric.
\end{proof}
This is a controlled spatial-design experiment with repeated readings
at each vertex, not a theorem about conditioning an arbitrary continuous
endpoint distribution on a point of probability zero. A random-endpoint
experiment requires a density or regression estimation step. On regular
$d$-dimensional meshes, $N_h\asymp h^{-d}$, the displayed bound has
order $h+\sqrt{\log(N_h/\zeta)/n}$. With a fixed total budget it gives
the familiar bias--variance balance, but no optimality for geometric
parameters or derivative loss is inferred from that calculation.
